% Folder 84, batch 07 — archivists' pages 121–140.
% Pass: Opus 5.5 (claude-opus-5-5), 2026-10-09 — first pass, unchecked against the pages by a human.
% Revised: Opus 5.5 (claude-opus-5-5), 2026-10-10 — re-read on the 700 dpi tiles; 44 \ill and 11 \uncertain resolved, 29 \ill and 21 \uncertain remain (new doubtful readings are offered as \uncertain).
%
% Even pages 122–140 are skipped: 122 is blank, 124–140 are typed pages
% (pp. 465–473) of a non-mathematical typescript used as scrap paper.
% Notation fixed for the batch: the scalar he writes as a long-tailed
% x / curly stroke is set as \chi throughout (it satisfies \sigma\chi = 2,
% p. 127, and "2 - \chi = 1", p. 139); the product of extensions is written
% * or \wedge as on the page, \vee where he uses it for the same object.
\documentclass[11pt,a4paper]{article}
\input{../preamble/grothendieck}

\folder{84}
\batch{7}
\pages{121}{140}
\foldertitle{[Formes quadratiques] : notes manuscrites (s.d.).}
\dating{[à partir de 1982-vers 1986]}
\watermark{Édition de démonstration}

\begin{document}

\section{Produit d'extensions et chi-scindages}

\page{121}
\note{la page reprend une construction commencée avant ce lot ; la notation $\mathfrak{X}$, $X_i$ y est déjà fixée}

Une \ill{} variante, pour $X_1, X_2, X_3$ fixés, la \uncertain{situation} universelle, où le \add{\uncertain{(universel)},} \ill{} sera.

$\mathfrak{X}_0 = X_1 \wedge X_2 \wedge X_3$, \struck{\ill{}} \uncertain{muni} du \uncertain{revêtement}
\[
  X_{1\,\mathfrak{X}_0},\quad X_{2\,\mathfrak{X}_0},\quad X_{3\,\mathfrak{X}_0}
  \qquad (X_{i\,\mathfrak{X}_0} = X_i \times_{\mathfrak{X}} \mathfrak{X}_0)
\]
\note{l'indice du produit fibré est lu \uncertain{$\mathfrak{X}$} ; il pourrait être $\mathfrak{X}_0$}

d'où $\mathfrak{X} \mapsto X_1 \times X_2 \times X_3$, …

\begin{tikzcd}[column sep=small, row sep=small, nodes={font=\scriptsize}]
  & X_{1\,\mathfrak{X}_0} \arrow[dr] & \\
  X_1 \times X_2 \times X_3 \arrow[ur, "{(\mathrm{pr}_1,\wedge)}"] \arrow[r, "{(\mathrm{pr}_2,\wedge)}"] \arrow[dr, "{(\mathrm{pr}_3,\wedge)}"'] & X_{2\,\mathfrak{X}_0} \arrow[r] & \mathfrak{X}_0 \\
  & X_{3\,\mathfrak{X}_0} \arrow[ur] &
\end{tikzcd}

ou \struck{\ill{}}
\[
  X_1 \times X_2 \times X_3 \longrightarrow \mathfrak{X}_0 = X_1 \wedge X_2 \wedge X_3
\]
est la multiplication
\[
  a_1, a_2, a_3 \longmapsto x_1 \wedge \alpha_2 \wedge x_3
\]
\note{la ligne porte bien $a_1, a_2, a_3$ à gauche et, à droite, ce qui se lit $x_1 \wedge \alpha_2 \wedge x_3$ ; l'incohérence des lettres est sur la page}

Une difficulté, c'est que la multiplication n'est pas \ill{} fini \uncertain{loc.\ libre}, ni $=$ plat (mais il \ill{} bien \ill{}).
\[
  B = A_1 \otimes_{O} A_2 \otimes_{O} A_3
  \;\rightleftarrows\;
  A_1 \otimes_{O} B_0 = B_1,
  \qquad B_0 = A_1 * A_2 * A_3
\]
\note{entre $B$ et $A_1 \otimes_O B_0$ deux flèches, marquées « 2(!) » ; au-dessus, un passage raturé d'où se détachent $\mathrm{Tr}_1$, $N_1$ et un « ? ». Sous $B$, la tour $O \xrightarrow{2} B_0 \xrightarrow{4(!)} B$, la seconde flèche annotée (\ill{} \ill{}), un arc marqué « 8 » de $O$ à $B$, et un trait de $B_0$ vers $A_1 \otimes_O B_0$ marqué « $\varepsilon 2$ »}

$\mathrm{Tr}_1$ doit être $B_0$-linéaire, $N_1$ $B_0$-\uncertain{quadratique} …

\page{123}

\begin{tikzcd}
  0 \arrow[r] & L \arrow[r] \arrow[d, "\alpha"] & E \arrow[r] \arrow[d, "\beta"] \arrow[l, bend right=40, "\pi"'] & M \arrow[r] \arrow[d, "\gamma"] \arrow[l, bend right=40, "\varpi"'] & 0 \\
  0 \arrow[r] & L' \arrow[r] & E' \arrow[r] \arrow[l, bend left=40, "\pi'"] & M' \arrow[r] \arrow[l, bend left=40, "\varpi'"] & 0
\end{tikzcd}

\[
  \alpha\,\pi(x) = \pi'\beta(x)
  \;\overset{?}{\Longleftrightarrow}\;
  \beta\varpi(y) = \varpi'\gamma(y)\ ?
\]
\[
  \text{\struck{$\beta$}}\ \pi(x) = \pi'\beta(x)
  \;\Longleftrightarrow\;
  \beta\varpi(x) \equiv
\]
\note{sous $\pi(x) = \pi'\beta(x)$, une accolade ; puis une ligne $\beta(\ldots - \varpi) = (\ldots\, \varpi')\beta$, à demi raturée, et une seconde, barrée, qui s'achève par $\beta$}

\begin{tikzcd}
  E_1 \arrow[d] & E_2 \arrow[d] \\
  E_1' & E_2'
\end{tikzcd}

$\Longrightarrow$

\begin{tikzcd}
  E_1 \wedge E_2 \arrow[r] & E_1'
\end{tikzcd}

\note{le but s'arrête sur $E'$, d'indice \uncertain{1} ; le second facteur n'est pas écrit}

\begin{tikzcd}[column sep=small, row sep=small, nodes={font=\scriptsize}]
  & E_1 \otimes E_2 \arrow[rr] & & E_1' \otimes E_2' \arrow[dr] & \\
  E_1 \otimes L_2,\ L_1 \otimes E_2 \arrow[ur] \arrow[dr, hook] & & & & E_1' \wedge E_2' \\
  & L_1 \otimes L_2 \arrow[rr] & & L_1' \otimes L_2' \arrow[ur] &
\end{tikzcd}

\note{la flèche $L_1' \otimes L_2' \to E_1' \wedge E_2'$ porte une annotation \ill{}}

\[
  x_1 \otimes x_2 \longmapsto \beta_1(x_1) \otimes \beta_2(x_2)
  \longrightarrow \beta_1(x_1) * \beta_2(x_2)
\]
\[
  = \pi_1'\beta_1(x_1) \otimes \alpha_2(\ldots)
\]
\note{la dernière ligne est surchargée : après $\pi_1'$ un symbole raturé, et le second facteur, écrit sur un $\gamma$ ou un $\beta$ raturé, porte un argument repassé, \uncertain{$\beta_2(x_2)$}}
\[
  \pi_1(x_1) \otimes u_2 \longrightarrow
  \underbrace{\alpha_1\pi_1(x_1)}_{\pi_1'\beta_1(x_1)} \otimes \alpha_2(x_2)
\]

$E_1 \wedge$
\note{la page s'arrête sur ce début de formule}

\page{125}
\[
  \mathcal{E}_{(p_1,q_1)(p_2,q_2)} = (M_1 \otimes M_2) \oplus (L_1 \otimes L_2)
  \xrightarrow{\ T_u\ }
  \mathcal{E}_{(p_1',q_1')(p_2',q_2')}
\]

\begin{tikzcd}
  \mathcal{E}_{(p_1,q_1)(p_2,q_2)} \arrow[r, "T_u"] \arrow[d, "T_{-\sigma b}"'] & \mathcal{E}_{(p_1',q_1')(p_2',q_2')} \arrow[d, "T_{-\sigma b'}"] \\
  \mathcal{E} \arrow[r, "T_{-u}"'] & \mathcal{E}
\end{tikzcd}

\note{au centre du carré, entouré : « commute ! »}
\[
  b' = b + \chi u
\]
\[
  u - \sigma b' = u - \sigma b - \sigma\chi u = -u - \sigma b
\]
\note{dans les deux lignes un signe a été surchargé ; le $+$ de $b' = b + \chi u$ est écrit sur un autre signe}

d'où un \add{can.}\ iso d'ext.\ $\chi$-scindées
\[
  E_1 \vee E_2 \simeq E_1 \wedge E_2
\]
\[
  E_1' \vee E_2' \;\simeq\; E_1' \wedge E_2'
\]
\note{cette seconde ligne est écrite entre les lignes, au-dessus de la phrase qui suit}

Je dis qu'il est fonctoriel pour les \emph{\uncertain{iso}}) : clair par construction. \add{Mais il y a aussi une notion} d'\emph{\uncertain{hom.}}\ d'ext.\ $\chi$-scindées. Ça marche \uncertain{aussi}

\drawing{deux carrés emboîtés, en perspective, formés des suites $L_i \to E_i \to M_i$, $L_i' \to E_i' \to M_i'$, $L_i'' \to E_i'' \to M_i''$ et $L_i''' \to E_i''' \to M_i'''$, reliés par des flèches obliques et des signes d'égalité ($M_i' \Vert M_i$) ; le dessin est trop surchargé pour être mis en \texttt{tikz-cd} sans choisir des flèches}

\begin{tikzcd}
  L_1 \otimes L_2 \arrow[r, hook] \arrow[d, Rightarrow] & E \arrow[r, two heads] & M_1 \otimes M_2 \arrow[ld, no head] \\
  L_1' \otimes L_2' & M_1' \otimes M_2' &
\end{tikzcd}

\note{$M_1' \otimes M_2'$ est écrit au-dessus de $M_1 \otimes M_2$, relié par un trait ; la position dans la grille est approximative}

\page{127}
\[
  \tilde b_i = \varpi_i - \chi q_i\,; \qquad
  b_i + \tilde b_i = \chi\,\mathrm{id} - \chi\,\mathrm{id} = 0
\]
\[
  \tilde b_i = -b_i
\]
\note{« $= -b_i$ » est doublement souligné ; le « $= 0$ » final de la seconde égalité est écrit petit, en bout de ligne}

$\tilde u_i$ tel que
\[
  q_i' = \underbrace{q_i}_{\mathrm{id} - p_i} - \tilde u_i
  = \text{\struck{$\chi$}}\,\mathrm{id} - \underbrace{(p_i + \tilde u_i)}_{p_i'} =
\]
\[
  \tilde u_i = -u_i
\]
\[
  \tilde u = \tilde b_1 \otimes \tilde u_2 + \tilde u_1 \otimes \tilde b_2 + \chi\, \tilde u_1 \otimes \tilde u_2
  = b_1 \otimes u_2 + u_1 \otimes b_2 + \chi\, u_1 \otimes u_2 = u
\]
\marginal{Mais alors $\tilde T_{\tilde u} = T_{-\tilde u} = T_{-u} = T_u^{-1}$}
\[
  \tilde b = \tilde b_1 \tilde b_2 = b_1 b_2
\]

\begin{tikzcd}
  \mathcal{E}_{(p_1,q_1)(p_2,q_2)} = M_1 \otimes M_2 \oplus L_1 \otimes L_2 \arrow[d, "\wr"', "T_{v(p_1,q_1)(p_2,q_2)}"] \\
  \mathcal{E}_{(\tilde p_1,\tilde q_1)(\tilde p_2,\tilde q_2)} = M_1 \otimes M_2 \oplus L_1 \otimes L_2
\end{tikzcd}

\[
  (= (\eta,\xi) \longmapsto (\eta,\ \xi + v(\eta))),
\]
\note{devant $v(\eta)$, un signe noirci, lu $+$}

$\chi$-scindage défini (comme \uncertain{g}) par
\[
  M_1 \otimes M_2 \xrightarrow{\ b = b_1 b_2\ } L_1 \otimes L_2
\]
$\chi$-sc.\ défini (comme \ill{}) par $-b$.

$T_v$ \struck{\ill{}}transforme le $\chi$-scindage défini par $b$ en le $\chi$-scindage défini \struck{par} $b + \chi v$ ; \ill{} \ill{} \ill{} que
\[
  b + \chi v = \underset{\displaystyle -b}{\overset{\sim}{b}}
  \quad\text{i.e.}\quad v \overset{!}{=} -\frac{2}{\chi}\, b
\]
Supposons que $\exists\, \sigma \in k$, tq.\ $\sigma\chi = 2$, \uncertain{alors} prendre
\[
  \boxed{v = -\sigma b}
\]

\page{129}
\note{feuillet écrit en largeur ; une barre verticale sépare une colonne de droite, transcrite après la colonne de gauche}
\[
  0 \to L_i \to E_i \to M_i \to 0, \qquad \pi_i : E_i \to L_i
\]
$\chi$-scindage ; \ill{}

$p_i : E_i \to L_i$ scindage, \quad $p = (p_1, p_2)$
\[
  \Theta_{p_i} : E_i \simeq M_i \oplus L_i
\]
\[
  b_i := \pi_i - \chi p_i : M_i \to L_i
\]
\[
  b = b_1 \otimes b_2 : M_1 \otimes M_2 \to L_1 \otimes L_2
\]
d'où $\chi$-trivialisation \ill{} \ill{} $\pi_{(p_1,p_2)}$ :
\[
  \mathcal{E}_{(p_1,p_2)} = (M_1 \otimes M_2) \oplus (L_1 \otimes L_2)
\]
\note{après $\mathcal{E}$, un indice raturé ; l'indice $(p_1,p_2)$ est récrit dessous}
\[
  \pi_{(p_1,p_2)} - \chi\, p = b_1 \otimes b_2
\]
\[
  \pi_{(p_1,p_2)} : (\eta,\xi) \longmapsto \text{\struck{\ill{}}}\ \chi\xi + b_1 \otimes b_2(\eta)
\]

\[
  p_i' = p_i - u_i, \quad u_i : M_i \to L_i \qquad (p' = (p_1', p_2'))
\]
d'où $b_1', b_2'$, avec
\[
  b_1' = b_1 + \chi u_1, \qquad b_2' = b_2 + \chi u_2
\]
\[
  \mathcal{E}_{p} \xrightarrow{\ \Theta_{p,p'}\ } \mathcal{E}_{p_1',p_2'},
  \qquad \Theta_{p,p'} = T_{-u} :
  (\eta,\xi) \longmapsto (\eta,\ \xi - u\eta)
\]
\note{au-dessus de la flèche, un symbole raturé ; sous $\Theta_{p,p'}$, une première écriture surchargée avant « $= T_{-u}$ », dont le $T$ et le signe sont noircis}
\[
  u = b_1 \otimes u_2 + u_1 \otimes b_2 + \chi\, u_1 \otimes u_2
\]
\emph{NB} $\Theta_{p,p'}$ transforme $\pi_{(p_1,p_2)}$
\[
  \pi_{p_1',p_2'} : (\eta,\xi) \longmapsto \text{\struck{\ill{}}}\ \chi\xi + (b_1' \otimes b_2')(\eta)
  = \pi_{(p_1,p_2)}(\eta,\xi) + u(\eta)
\]

\page{131}
\[
  \text{\struck{$\Theta_{p+u}(x) =$}}
  \qquad
  \Theta_p(x) = \psi(x) \oplus p(x)
\]
\[
  \Theta_{p+u}(x) = \psi(x) \oplus \bigl(p(x) + u\,\psi(x)\bigr)
\]
\[
  \boxed{b_i = (\pi_i - \chi p_i)}
\]

\begin{tikzcd}[column sep=small, row sep=small, nodes={font=\scriptsize}]
  & & E_1 \vee E_2 \arrow[dll, "\Theta_{p_1 * p_2}"', "\sim"] \arrow[d, "\Theta_{p_1' * p_2'}", "\wr"'] \arrow[drr, "\Theta_{p_1'' * p_2''}", "\sim"'] & & \\
  M \oplus L \arrow[rr, "{\Theta_{(p_1,p_2)(p_1',p_2')}}", "{(u_1,u_2,u)}"'] & & M \oplus L \arrow[rr, "{\Theta_{(p_1',p_2')(p_1'',p_2'')}}", "{(u_1',u_2',u')}"'] & & M \oplus L
\end{tikzcd}

\[
  \mathrm{Hom}(M_i, L_i) \ni b_i,
  \qquad
  (\eta,\xi) \longmapsto (\eta,\ \xi - u\eta)
\]
\[
  \mathrm{Hom}(M_i, L_i) \ni b_i' = b_i + \chi u_i,
  \qquad
  (\eta,\xi) \longmapsto (\eta,\ \xi - u'\eta)
\]
\[
  b_i'' = b_i' - \chi u_i'
\]
\note{un arc, étiqueté $u_1'', u_2'', u''$ et $\Theta_{(p_1,p_2)(p_1'',p_2'')}$, joint le premier $M \oplus L$ au troisième : $(\eta,\xi) \to (\eta, \xi + \ldots)$}
\[
  b_1 \otimes b_2 \in \mathrm{Hom}(M, L),
  \qquad
  (\pi - \chi\, p_1 * p_2)^{\Theta_{p_1 * p_2}}
\]
\[
  b_1' \otimes b_2' = (b_1 - \chi u_1) \otimes (b_2 - \chi u_2) = b_1 \otimes b_2 - \chi u
\]
\[
  b_1'' \otimes b_2'' = b_1' \otimes b_2' - \chi u' = b_1 \otimes b_2 - \chi\underbrace{(u + u')}_{u''}
\]
\note{les signes de $b_i' = b_i \pm \chi u_i$ varient d'une ligne à l'autre ; ils sont transcrits tels quels. Dans $= b_1 \otimes b_2 - \chi u$, le signe est surchargé, un $+$ et un $-$ l'un sur l'autre}
\[
  \left\{
  \begin{aligned}
    p_1' * p_2' &= p_1 * p_2 - u \\
    p_1'' * p_2'' &= p_1' * p_2' - u' = p_1 * p_2 - u'' \\
    &\Longrightarrow u'' = u + u'
  \end{aligned}
  \right.
\]

\[
  \underbrace{b_1, b_2;\ u_1, u_2}_{\mathrm{I}}
  \qquad
  \underbrace{u_1',\ u_2'}_{\mathrm{I}}
\]
\[
  \begin{array}{c}
    (p_1,p_2) \\ \downarrow \\ (p_1',p_2') \\ \downarrow \\ (p_1'',p_2'') \\[2pt] \hline
    (p_1,p_2) \\ \downarrow \\ (p_1'',p_2'')
  \end{array}
  \quad
  \left\{
  \begin{aligned}
    u &= b_1 \otimes u_2 + u_1 \otimes b_2 + \chi\, u_1 \otimes u_2 \\
    u' &= b_1' \otimes u_2' + u_1' \otimes b_2' + \chi\, u_1' \otimes u_2' \\
    u'' &= b_1 \otimes u_2'' + u_1'' \otimes b_2 + \chi\, u_1'' \otimes u_2'' \\
    &\quad (\text{si } u_1'' = u_1 + u_1',\ u_2'' = u_2 + u_2')
  \end{aligned}
  \right.
\]
\[
  \left\{
  \begin{aligned}
    b_1' &= b_1 + \chi u_1, & b_2' &= b_2 + \chi u_2 \\
    b_1'' &= b_1' + \chi u_1', & b_2'' &= b_2 + \chi u_2'
  \end{aligned}
  \right.
  \qquad
  \begin{aligned}
    b &= b_1 \otimes b_2 \\
    b' &= b_1' \otimes b_2' \\
    b'' &= b_1'' \otimes b_2''
  \end{aligned}
\]
\note{$b_2'' = b_2 + \chi u_2'$ (et non $b_2' + \ldots$) est sur la page. Une ligne barrée suit le système : $= (b_1 - \chi u_1) \otimes u_1' + u_2' \otimes (b_2 - \chi u_2) + \chi u_1' \otimes \ldots$}

1°) Je dis que $u'' = u + u'$, \struck{$b_1'' = b_1' + \chi u_1'',\ b_2'' =$}
\[
  u'' = b_1 \otimes (u_2 + u_2') + (u_1 + u_1') \otimes b_2
  + \chi\, (u_1 + u_1') \otimes (u_2 + u_2')
\]
\[
  = u + b_1 \otimes u_2' + \underbrace{u_1' \otimes b_2 + \chi\, u_1' \otimes u_2}_{u_1' \otimes (b_2 + \chi u_2)}
  + \chi\, u_1 \otimes u_2' + \chi\, u_1' \otimes u_2'
\]
\[
  \ldots + (b_1 + \chi u_1) \otimes u_2' + \chi\, u_1' \otimes u_2'
\]
\note{à droite, un « $=$ » suivi d'un symbole raturé ; le calcul s'arrête là}

\page{133}

\begin{tikzcd}[column sep=small, nodes={font=\scriptsize}]
  & E_1 \vee E_2 \arrow[dl, "\Theta_{p_1 * p_2}"'] \arrow[dr, "\Theta_{p_1' * p_2'}"] & \\
  (M_1 \otimes M_2) \oplus (L_1 \otimes L_2) \arrow[rr, "\Lambda_u"'] & & (M_1 \otimes M_2) \oplus (L_1 \otimes L_2)
\end{tikzcd}

\[
  \Lambda_u : (\eta,\xi) \longmapsto (\eta,\ \xi + u(\eta)),
  \qquad b_1, b_2 \quad\leadsto\quad b_1', b_2'
\]
\[
  \Theta_{p_1' * p_2'} = \Theta_{p_1 * p_2}\ \ill{}\ u
\]
\[
  u = b_1(p_1) \otimes u_2 + u_1 \otimes b_2(p_2) + \chi\, u_1 \otimes u_2
\]
\[
  b_1' = b_1 - \chi u_1, \qquad b_2' = b_2 - \chi u_2
\]
\[
  p_1' = p_1 - u_1\psi_1, \qquad p_2' = p_2 - u_2\psi_2,
  \qquad u_1 : M_1 \to L_1,\quad u_2 : M_2 \to L_2
\]
\note{le signe devant $u_i\psi_i$ est surchargé, lu $-$}
\[
  p_1' * p_2' = p_1 * p_2 - \Bigl[\underbrace{(\pi_1 - \chi p_1)}_{b_1(p_1)} \otimes u_2
  + u_1 \otimes \underbrace{(\pi_2 - \chi p_2)}_{b_2(p_2)}\Bigr]
  - \chi\, u_1 \otimes u_2
\]
\[
  u_1 * u_2 = -\chi\, u_1 \otimes u_2
\]
\marginal{une flèche montre le $-$ devant $\chi\, u_1 \otimes u_2$ : signe étrange, \ill{} \ill{} \struck{\ill{}}}

Ces signes sont dûs au signe $-\mathrm{id}$ dans le diag.\ \ill{} \uncertain{commutatif} \ill{} pour une \uncertain{catégorie} \struck{\ill{}} de $\underline{O}$ par $\mathcal{H}$

\begin{tikzcd}
  0 \arrow[r] & \mathcal{H} \arrow[r, "\alpha"] \arrow[d, "-\mathrm{id}"'] & \mathcal{E} \arrow[r, "\psi"] \arrow[d, "\Phi"] & \underline{O} \arrow[r] \arrow[d, no head, "\Vert" description] & 0 \\
  0 \arrow[r] & \mathcal{H} \arrow[r] & \mathcal{H}om^{!}(\mathcal{E}, \mathcal{H}) \arrow[r] & \underline{O} \arrow[r] & 0
\end{tikzcd}

\[
  \Phi(x)(y) = -x\,\psi(y) + y\,\psi(x)
\]
\[
  \mathcal{H}om(\underline{O}, \mathcal{H})
  = \text{\struck{$\mathcal{H}om \ldots f \mapsto \lambda \ldots \exists \lambda$}}
  \ \bigl\{ f : \mathcal{E} \to \mathcal{H} \bigm| f \circ \alpha = \alpha(\lambda)\,\mathrm{id}_{\mathcal{H}} \bigr\}
\]
\[
  \xi_1,\ \eta_1 \longmapsto \chi\eta_1 + b_1(\xi_1),
  \qquad \xi_1 \in M_1 \otimes M_2,\ \eta_1 \in L_1 \otimes L_2
\]
\note{au-dessus de cette ligne, une seconde flèche $\longleftrightarrow$ vers un symbole raturé}

Si on \uncertain{prend}
\[
  p_1' = p_1 - u_1', \qquad p_2' = p_2 - u_2', \qquad p_1' * p_2' = p_1 * p_2 - u'
\]
d'où
\[
  \left\{
  \begin{aligned}
    u_1' &= -u_1 \\ u_2' &= -u_2 \\ u_3' &= -u_3
  \end{aligned}
  \right.
  \quad\leadsto\ldots\quad
  u' = -b_1 \otimes u_2 - u_1 \otimes b_2 + \chi\, u_1 \otimes u_2
\]
\[
  \boxed{u' = b_1 \otimes u_2' + u_1' \otimes b_2 + \chi\, u_1' \otimes u_2'}
\]
\marginal{\struck{\ill{}} la formule \uncertain{transitive}…}

\page{135}

\begin{tikzcd}
  0 \arrow[r] & L \arrow[r] \arrow[d, "-\mathrm{id}_L"'] & E \arrow[r, "\psi"] \arrow[d, "\wr"] & \underline{O} \arrow[r] \arrow[d, no head, "\Vert" description] & 0 \\
  0 \arrow[r] & L \arrow[r] & \mathcal{H}om^{!}(E, L) \arrow[r] & \underline{O} \arrow[r] & 0
\end{tikzcd}

\[
  x \longmapsto \bigl[\, y \longmapsto (-x\,\psi(y) + y\,\psi(x)) \,\bigr]
\]
\[
  E \xrightarrow{\ \psi(x)\ } L
\]
\note{l'étiquette $\psi(x)$ est lue sous la flèche ; la flèche n'est pas autrement commentée}
\[
  {}_{\lambda_1}E_1 * {}_{\lambda_2}E_2 \simeq
\]
\[
  \boxed{
  \begin{aligned}
    (f_1 * f_2)(x_1 * x_2) &= \psi_1(f_1)\,\pi_1(x_1) \otimes f_2(x_2)
    + f_1(x_2) \otimes \psi_2(f_2)\,\pi_2(x_2) \\
    &\quad - \chi\, f_1(x_1) \otimes f_2(x_2) \\
    (f_1 * f_2)(u_1 \otimes u_2) &= \lambda_1\lambda_2\, u_1 \otimes u_2
  \end{aligned}}
\]
\note{$f_1(x_2)$ dans le deuxième terme est sur la page ; devant $\pi_1$ un premier symbole est raturé}
\[
  (f_1 * u_2)(x_1 * x_2) = \psi_1(f_1)\,\pi_1(x_1) \otimes \underbrace{u_2(x_2)}_{u_2(\psi_2(x_2))}
  - \chi\, f_1(x_1) \otimes u_2(x_2)
\]
\[
  \text{\struck{$(u_1 * f_2)(x_1 * x_2)$}}
\]
i.e.
\[
  (f_1 * u_2)(x_1 * x_2) = \bigl(\underbrace{\psi_1(f_1)\,\pi_1 - \overset{\psi_1(\pi_1)}{\chi}\, f_1}_{\text{nul sur } L_1}\bigr) \otimes u_2\,(x_1 \otimes x_2)
\]
\[
  \bigl( [\pi_1', f_1']' \otimes u_2 \bigr)(x_1 \otimes x_2)
\]

d'où
\[
  \left\{
  \begin{aligned}
    (u_1 * f_2)(x_1 * x_2) &= \bigl(u_1 \otimes [\pi_2', f_2']'\bigr)(x_1 \otimes x_2) \\
    (u_1 * u_2)(x_1 * x_2) &= -\chi\, u_1(x_1) \otimes u_2(x_2)
  \end{aligned}
  \right.
\]
\marginal{une flèche montre le $-$ : signe étrange !}

\[
  (x_1 + u_1) * (x_2 + u_2) = x_1 * x_2
  + \underbrace{x_1 * u_2}_{\underbrace{\pi_1(x_1)}_{b_1} \otimes u_2}
  + \underbrace{u_1 * x_2}_{u_1 \otimes \underbrace{\pi_2(x_2)}_{b_2}}
  + \underbrace{u_1 * u_2}_{\chi\, u_1 \otimes u_2}
\]

\page{137}
\[
  \text{\struck{$(f_1 * f_2)(x_1 * x_2) =$}}
  \qquad
  (f_1 * f_2) \mid E_1 \otimes E_2 : E_1 \otimes E_2 \to L_1 \otimes L_2
\]
\[
  \boxed{
  \begin{aligned}
    f_1 * f_2\,(x_1 * x_2) &= \lambda_1\pi_1(x_1) \otimes f_2(x_2)
    + f_1(x_1) \otimes \lambda_2\pi_2(x_2) - \chi\, f_1(x_1) \otimes f_2(x_2)
  \end{aligned}}
\]
i.e.
\[
  f_1 * f_2 \mid E_1 \otimes E_2 \ \bigl[E_1 \otimes E_2 \to L_1 \otimes L_2\bigr]
  \doteq \lambda_1\pi_1 \otimes f_2 + f_1 \otimes \lambda_2\pi_2 - \chi\, f_1 \otimes f_2
\]
\[
  f_1 * f_2\,(u_1 \otimes u_2) = \lambda_1\lambda_2\,(u_1 \otimes u_2)
\]
$f_1 \longrightarrow x$\note{isolé, à gauche de « Vérifions »}\quad Vérifions
\[
  f_1 * f_2\,(x_1 * u_2) = \lambda_1\lambda_2\, \pi_1(x_1) \otimes u_2
\]
\[
  \lambda_1\pi_1(x_1) \otimes \lambda_2 u_2
  + \text{\struck{$f_1(x_1) \otimes \lambda_2\chi\, u_2$}}
  - \chi\, \underbrace{f_1(x_1) \otimes \lambda_2 u_2}
\]
\marginal{OK}
\[
  (\pi_1 * \pi_2)(x_1 * x_2) = \chi\,\{\pi_1(x_1) \otimes \pi_2(x_2)\}
\]
\[
  (\pi_1 * \pi_2)(u_1 \otimes u_2) = \chi^2\, u_1 \otimes u_2
  \qquad
  \frac{1}{\chi}\,\pi_1 \otimes \pi_2 = \pi
\]
\[
  x_1 * x_2 \equiv \text{\struck{$\psi_1(x_1)$}}
\]
\note{à droite d'une barre verticale :}
\[
  \begin{array}{ll}
    L_i & p_i \\ M_i & q_i
  \end{array}
  \qquad
  1 + \underbrace{p_1 q_1},\quad 1 + \underbrace{p_2 q_2}
  \qquad 1 + (p_1 p_2)(q_1 q_2)
\]
\[
  \text{\struck{$[x,y] =$}}\quad
  x \longrightarrow \text{\struck{$\psi(x) = \ldots$}}\ \psi(x)\,T - \chi
\]
\note{le dernier terme, après $\chi$, est raturé}

\page{139}
\[
  \begin{array}{l}
    M_1 \\ \downarrow\ \ q_1 \\ E_1 \\ \alpha_1 \downarrow\ \ p_1 \\ L_1
  \end{array}
\]
\note{$q_1$ et $p_1$ sont écrits comme des flèches courbes remontant de $M_1$ vers $E_1$ et de $E_1$ vers $L_1$ ; le sens est celui de la page}

\begin{tikzcd}[column sep=small, row sep=small, nodes={font=\scriptsize}]
  (E_1 \otimes L_2) \oplus (L_1 \otimes E_2) \arrow[r, Rightarrow] \arrow[dr, "\pi"'] & E_1 \otimes E_2 \arrow[r] \arrow[d, "x_1 \otimes x_2 \mapsto x_1 * x_2"] & M_1 \otimes M_2 \arrow[r] \arrow[d, no head, "\wr"] & 0 \\
  & L_1 \otimes L_2 \arrow[r] & E \arrow[r] \arrow[l, bend left=40, "p"] & M_1 \otimes M_2 \arrow[l, bend left=30, "{m_1 \otimes m_2 \mapsto x_1 * x_2}"]
\end{tikzcd}

\note{la ligne du haut est surmontée d'un $M_1 \otimes M_2$ raturé ; au centre du carré, le mot \uncertain{cocart.}, et sur la flèche oblique de $E_1 \otimes E_2$ vers $E$ une étiquette \uncertain{$\lambda_1$} ; la flèche $p$ est lue sur une lettre courbe, \uncertain{$p$}}
\[
  q(m_1 \otimes m_2) = q(m_1) * q(m_2)
\]
\[
  p(x_1 * x_2) = x_1 * x_2 - q(m_1 \otimes m_2)
\]
\[
  = x_1 * x_2 - q_1(m_1) * q_2(m_2)
\]
\[
  q(m_1) = x_1 - p_1(x_1), \qquad q(m_2) = x_2 - p_2(x_2)
\]
\[
  p(x_1 * x_2) = x_1 * x_2 - \bigl(x_1 - p_1(x_1)\bigr) * \bigl(x_2 - p_2(x_2)\bigr)
\]
\[
  = x_1 * p_2(x_2) + p_1(x_1) * x_2 - \underline{p_1(x_1) * p_2(x_2)}
\]
\[
  \text{\struck{$= \chi\,(x_1 \otimes p_2(x_2) + p_1(x_1) \otimes x_2$}}
\]
\[
  \text{\struck{$-\ \chi\, p_1(x_1) \otimes p_2(x_2))$}}
\]
\[
  u_1 * u_2 = \chi\, u_1 \otimes u_2,
  \qquad
  \text{\struck{$p(x_1 *$}}\ \ p(u_1 * u_2) = \chi\, u_1 \otimes u_2
\]
\[
  \boxed{
  (f_1 * f_2)(x_1 * x_2) = \lambda_1 x_1 * f_2(x_2) + f_1(x_1) * \lambda_2 x_2
  - \chi\, f_1(x_1) \otimes f_2(x_2)}
\]
\note{sous $\lambda_1 x_1 * f_2(x_2)$, un renvoi $\Vert$ vers $\lambda_1 \pi_1(x_1)$}

\struck{si $x_1, x_2 \in L$}

si $x_1 = u_1 \in L_1$, $x_2 = u_2 \in L_2$
\[
  2\lambda_1 \ldots\, \chi\,(x_1 \otimes x_2) - \chi\, \lambda_1\lambda_2\, x_1 \otimes x_2
\]
\[
  = \text{\struck{$(2 - \chi)\,\lambda_1\lambda_2\, u_1 \otimes u_2$}}
  \qquad 2 - \chi = 1
\]
\note{la première ligne est surchargée : entre $2\lambda_1$ et $(x_1 \otimes x_2)$, un $\lambda_2$ incertain est écrit sur une rature, et le $\chi$ du second terme est lui aussi repassé}
\[
  (f_1 * f_2)(u_1 \otimes u_2) = \lambda_1\lambda_2\, u_1 \otimes u_2
  = f_1(u_1) \otimes f_2(u_2)
\]
\[
  \underbrace{xy - (x - x)(y - y)}_{\displaystyle xy + yx - xy} = xy
\]
\note{un grand trait oblique et une petite figure (une lettre ornée, \ill{}) séparent ce calcul du reste de la page}

\end{document}
